\chapter{Comparaison des méthodes sur la campagne}
\section{Bilan des exécutions et des garanties}
Les 168 exécutions des trois méthodes approchées fournissent une solution
réalisable. La PLNE est tentée sur quatre cas : \code{ulysses16} et
\code{ulysses22}, avec $p=3$ et $p=5$. Les 20 autres cas sont ignorés selon
les seuils configurés. Le tableau~\ref{tab:exacts} présente toutes les
résolutions exactes effectivement lancées.

\begin{table}[htbp]\centering
{\small\setlength{\tabcolsep}{3.5pt}\begin{tabular}{lrrrrrl}
\toprule
Instance & $p$ & Coût & Borne inf. & Écart$^*$ & Temps (s) & Optimum \\
\midrule
ulysses16 & 3 & 68{,}2731 & 68{,}2731 & 0{,}00\,\% & 0{,}089 & Prouvé \\
ulysses16 & 5 & 65{,}3775 & 65{,}3775 & 0{,}00\,\% & 1{,}772 & Prouvé \\
ulysses22 & 3 & 83{,}1382 & 83{,}1382 & 0{,}00\,\% & 0{,}489 & Prouvé \\
ulysses22 & 5 & 71{,}5707 & 66{,}3780 & 7{,}26\,\% & 5{,}018 & Non prouvé \\
\bottomrule
\end{tabular}
}
\caption{Solutions PLNE, bornes et statut de preuve, sous une limite CBC de cinq secondes.}
\label{tab:exacts}
\sourcecalcul{solutions JSON et journaux CBC enregistrés. L'écart$^*$ est
$100(C-B)/|C|$, avec la borne issue du journal ; les valeurs sont arrondies pour l'affichage.}
\end{table}

Les deux cas à trois stations sont résolus avec une preuve. Sur
\code{ulysses16} à cinq stations, CBC prouve également l'optimalité de la
valeur 65,3775. Sur \code{ulysses22} à cinq stations, il fournit une solution
de coût 71,5707, mais sa borne inférieure reste à 66,378 environ.
L'écart sauvegardé vaut 7,26\,\% selon la convention du projet. L'égalité
éventuelle d'une heuristique avec 71,5707 ne permet donc pas d'annoncer
un optimum sur ce cas.

Les seuils de taille ont un effet méthodologique : le rapport ne mesure pas
la capacité de CBC à résoudre les six grandes instances. Il ne peut donc
conclure que ces instances seraient insolubles par PLNE. Il constate seulement
qu'elles n'ont pas été lancées dans cette campagne de démonstration.

\section{Gains face à la construction}
Les deux recherches améliorent le coût moyen de la construction dans chacun
des 24 cas. Le tableau~\ref{tab:gains} résume les gains par cas, puis la
figure~\ref{fig:gains} permet d'en lire les différences selon l'instance et $p$.
\begin{table}[htbp]\centering
\begin{tabular}{lrrr}
\toprule
Méthode & Minimum & Médiane & Maximum \\
\midrule
Recherche locale & 0{,}80\,\% & 11{,}79\,\% & 39{,}74\,\% \\
Recherche tabou & 2{,}63\,\% & 13{,}67\,\% & 39{,}74\,\% \\
\bottomrule
\end{tabular}

\caption{Distribution descriptive des 24 gains moyens face à la construction, à $\alpha=1$.}
\label{tab:gains}\sourcetable
\end{table}

Les gains varient fortement : une construction déjà proche d'une bonne
sélection laisse peu de marge, alors qu'une sélection initiale mal adaptée
peut être beaucoup améliorée. Cette amplitude rappelle que la qualité de la
référence influence le gain relatif. Un grand pourcentage ne suffit pas à
qualifier une méthode de performante face à toutes les références possibles.

La moyenne tabou est meilleure que la moyenne locale dans 22 cas, et identique
à la précision considérée dans deux cas. Aucune moyenne locale ne bat la
moyenne tabou dans ce protocole. Cela décrit les réglages testés : 100
évaluations par recherche locale et 1\,500 par recherche tabou. Le résultat
ne sépare pas l'avantage de la mémoire tabou de celui d'un échantillon plus
large et d'un plus grand nombre de reconstructions du cycle.

\figureRS[0.93]{gains_campagne.pdf}{Gains moyens des recherches face à la construction.
Chaque case compare une instance et une valeur de $p$, avec $\alpha=1$.}
{fig:gains}
\FloatBarrier

\section{Un cas avec optimum prouvé : Ulysses16, cinq stations}
Le tableau~\ref{tab:u16} donne une comparaison où le coût optimal est connu
dans la métrique retenue par le projet. La construction vaut 89,1917.
La recherche locale atteint une moyenne de 65,5512 avec un écart-type de
0,1505. La recherche tabou atteint 65,3775 sur chacune des trois graines,
valeur égale à l'optimum prouvé par CBC.
\begin{table}[htbp]\centering
{\small\setlength{\tabcolsep}{4pt}\begin{tabular}{lrrrr}
\toprule
Méthode & Coût moyen & $s$ & Écart à l'optimum & Temps (ms) \\
\midrule
Construction & 89{,}1917 & \textemdash & 36{,}426\,\% & 0{,}16 \\
Locale & 65{,}5512 & 0{,}1505 & 0{,}266\,\% & 2{,}49 \\
Tabou & 65{,}3775 & 0{,}0000 & 0{,}000\,\% & 42{,}28 \\
PLNE / CBC & 65{,}3775 & \textemdash & 0{,}000\,\% & 1\,772{,}19 \\
\bottomrule
\end{tabular}
}
\caption{Ulysses16, $p=5$, $\alpha=1$ : coût, dispersion et écart à l'optimum prouvé.}
\label{tab:u16}\sourcetable
\end{table}

La dispersion nulle du tabou signifie seulement que les trois coûts observés
coïncident. Elle ne prouve pas qu'une autre graine ou un autre problème aurait
le même comportement. Ici, l'égalité de coût avec une valeur certifiée permet
cependant de conclure que les solutions tabou observées sont optimales pour
ce cas, sous les tolérances numériques de comparaison.

L'écart à l'optimum est renseigné uniquement sur les trois cas certifiés du
tableau~\ref{tab:gap}. Les autres cases de l'étude ne reçoivent pas d'écart
à un optimum non établi.
\begin{table}[htbp]\centering
\begin{tabular}{lrrrr}
\toprule
Instance & $p$ & Construction & Locale & Tabou \\
\midrule
ulysses16 & 3 & 65{,}961\,\% & 0{,}000\,\% & 0{,}000\,\% \\
ulysses16 & 5 & 36{,}426\,\% & 0{,}266\,\% & 0{,}000\,\% \\
ulysses22 & 3 & 65{,}653\,\% & 0{,}000\,\% & 0{,}000\,\% \\
\bottomrule
\end{tabular}

\caption{Écarts de coût moyen aux optima prouvés disponibles.}
\label{tab:gap}\sourcetable
\end{table}

\section{Qualité et durée : un compromis visible}
Les constructions sont très rapides et les améliorations demandent davantage
de temps. Le tableau~\ref{tab:temps} fixe $p=5$ pour faciliter la lecture des
ordres de grandeur. Le tabou met davantage de temps que la recherche locale,
ce qui est cohérent avec le nombre de candidats évalués. La différence de
durée n'est pas strictement proportionnelle à 15 : les ensembles de stations
et le nombre de changements 2-opt diffèrent entre trajectoires.

\begin{table}[htbp]\centering
\begin{tabular}{lrrr}
\toprule
Instance & Construction (ms) & Locale (ms) & Tabou (ms) \\
\midrule
ulysses16 & 0{,}16 & 2{,}49 & 42{,}28 \\
ulysses22 & 0{,}20 & 2{,}91 & 34{,}20 \\
att48 & 0{,}58 & 5{,}04 & 62{,}18 \\
eil51 & 0{,}66 & 5{,}29 & 67{,}04 \\
berlin52 & 0{,}66 & 5{,}91 & 74{,}36 \\
st70 & 1{,}04 & 7{,}26 & 89{,}81 \\
kroA100 & 1{,}87 & 10{,}27 & 118{,}40 \\
rd100 & 1{,}88 & 10{,}51 & 120{,}93 \\
\bottomrule
\end{tabular}

\caption{Temps moyens observés à $p=5$ et $\alpha=1$, en millisecondes ; une exécution
pour la construction, trois graines pour les recherches.}
\label{tab:temps}\sourcetable
\end{table}

Sur \code{kroA100}, à cinq stations, la moyenne de coût passe de 68\,427,57
pour la construction à 60\,284,12 pour la locale et 58\,433,18 pour le tabou.
La dispersion est aussi différente : les écarts-types observés sont
respectivement 1\,653,97 et 59,57 pour les deux recherches. Le tabou est ici
plus régulier sur trois graines, en contrepartie d'un temps moyen d'environ
118,40 ms, contre 10,27 ms pour la locale.

\section{Lire l'effet du nombre de stations sans le confondre avec la taille}
Les figures~\ref{fig:att48-camp} et~\ref{fig:kro-camp} conservent la même
instance et font varier $p$. Elles permettent de relier la valeur de l'objectif,
le temps et le gain face à la construction. Elles ne comparent pas directement
le coût brut de \code{att48} à celui de \code{kroA100}.

La cardinalité modifie simultanément le nombre de stations à choisir, la
taille du cycle et le voisinage d'échange. Le nombre de couples
$(p-1)(n-p)$ n'est pas constant lorsque $p$ varie. Même avec le même budget
de candidats, le degré d'exploration du voisinage change. Une interprétation
du seul coût final doit conserver ces différences à l'esprit.

\bilan{Les recherches réduisent les coûts de la construction dans les 24 cas.
Le tabou obtient de meilleurs coûts moyens dans ce protocole, avec un effort
et un temps supérieurs. Les garanties d'optimalité restent limitées aux cas
pour lesquels une preuve a effectivement été obtenue.}

\figureRS[0.96]{comparaison_att48.pdf}{Comparaison par cardinalité sur att48 :
coût moyen, temps moyen et gain relatif face à la construction.}{fig:att48-camp}
\figureRS[0.96]{comparaison_kroA100.pdf}{Même lecture sur kroA100. Les solutions exactes
ne sont pas lancées sur cette instance dans la campagne.}{fig:kro-camp}
\FloatBarrier
